\documentclass[11pt]{article}
\usepackage{amsmath,amssymb,amsthm,geometry,hyperref}
\geometry{margin=1in}
\newtheorem{theorem}{Theorem}
\newtheorem{lemma}{Lemma}
\newcommand{\Tr}{\operatorname{Tr}}
\begin{document}
\title{A rank-one-complement obstruction to anticoherence:\\seventeen qubits cannot be sixth-order AC}
\author{Exact catalogue calculation}
\date{3 November 2026}
\maketitle
\begin{theorem}\label{thm}
For every integer $r\ge1$, there is no pure $2r$-anticoherent state on
$N=(r+1)^2+1$ symmetric qubits. In particular,
\[
 \boxed{K((r+1)^2+1,2r)=0,\qquad K(17,6)=0.}
\]
The conclusion applies to arbitrary complex amplitudes.
\end{theorem}
This is an analytical representation-theoretic proof. The accompanying
\texttt{verify\_channel.py} checks the physical Clebsch--Gordan normalization,
operator product ranks, and covariant channel for the seventeen-qubit case
with exact algebraic arithmetic. Numerical optimization is not used.

\section*{1. Rectangular spin tensors and their orthogonal images}
Write $V_{N/2}=\operatorname{Sym}^N(\mathbb C^2)$, of dimension $N+1$.
A pure state $\rho=|\psi\rangle\langle\psi|$ is $t$-AC precisely when
\[
 \Tr(\rho T_{Lq})=0\qquad(1\le L\le t,\ -L\le q\le L)
\]
for all irreducible spin tensors on its spin space. Consequently,
\begin{equation}\label{mom}
 \langle\psi|A|\psi\rangle=\frac{\Tr A}{N+1}
\end{equation}
whenever $A$ has only tensor ranks $0,\ldots,t$.

Consider Hilbert--Schmidt orthonormal tensor maps
\[
 E_{\ell q}:V_{N/2}\longrightarrow V_{(N-2)/2},
 \qquad \ell=1,\ldots,r,\quad q=-\ell,\ldots,\ell.
\]
All these irreducible representations occur in
$\operatorname{Hom}(V_{N/2},V_{(N-2)/2})$, whose tensor ranks are
$1,2,\ldots,N-1$; the sizes used below satisfy $r\le N-1$.
There are
\[
 M=\sum_{\ell=1}^{r}(2\ell+1)=r(r+2)=(r+1)^2-1
\]
such maps. Products $E_{\ell q}^{\dagger}E_{\ell' q'}$ have ranks at most
$\ell+\ell'\le2r$. Thus if $\psi$ is $2r$-AC, \eqref{mom} and
Hilbert--Schmidt orthogonality give the exact Gram identity
\begin{equation}\label{gram}
 \langle E_{\ell q}\psi,E_{\ell' q'}\psi\rangle
 =\frac{\delta_{\ell\ell'}\delta_{qq'}}{N+1}.
\end{equation}
In particular $M$ nonzero orthogonal vectors fit into the output space of
dimension $N-1$, giving the familiar necessary bound
\begin{equation}\label{bound}
 \boxed{N\ge(r+1)^2}\qquad\text{for a pure $2r$-AC state.}
\end{equation}

\section*{2. Covariance transfers anticoherence to a rank-one complement}
Define the completely positive map
\[
 \Phi(A)=\sum_{\ell=1}^{r}\sum_{q=-\ell}^{\ell}
 E_{\ell q}A E_{\ell q}^{\dagger}.
\]
For each $\ell$ the Kraus operators contain the entire tensor multiplet.
Under a rotation their $q$ indices mix by a unitary spin-$\ell$ matrix,
so $\Phi$ is covariant:
\[
 \Phi(D_{N/2}(g)A D_{N/2}(g)^{\dagger})
 =D_{(N-2)/2}(g)\Phi(A)D_{(N-2)/2}(g)^{\dagger}.
\]
The operator spaces of both irreducible spin representations contain each
allowed tensor rank once. Schur's lemma therefore implies that $\Phi$ and
its Hilbert--Schmidt adjoint preserve tensor rank, possibly annihilating
it. In particular, if $\rho$ has no ranks $1,\ldots,2r$, neither does
$\Phi(\rho)$. This implication requires no nonvanishing channel multiplier.

By \eqref{gram},
\[
 Q=(N+1)\Phi(\rho)
\]
is an orthogonal projector of rank $M$. If its complement has rank one,
there is a normalized output vector $\phi$ such that
\[
 I_{N-1}-Q=|\phi\rangle\langle\phi|.
\]
The identity has only scalar tensor rank, while $Q$ has no ranks
$1,\ldots,2r$. Hence $\phi$ is a pure $2r$-AC state on $N-2$ qubits.
This is the rank-one-complement transfer used in the theorem.

\section*{3. The forbidden sizes}
Put $D=(r+1)^2$ and suppose $N=D+1$. Then
\[
 \dim V_{(N-2)/2}=N-1=D,\qquad\operatorname{rank}Q=M=D-1.
\]
The complement is therefore one-dimensional and supplies a pure $2r$-AC
state on $N-2=D-1$ qubits. That violates \eqref{bound}, which requires at
least $D$ qubits. This contradiction proves Theorem~\ref{thm}.
The first instances are
\[
 (N,t)=(5,2),(10,4),(17,6),(26,8),(37,10),\ldots.
\]
The theorem excludes these specific sizes; it is not an existence claim
at any neighboring size.

\section*{4. The seventeen-qubit case explicitly}
For $N=17$ and $r=3$, the fifteen tensors of ranks one, two and three map
spin $17/2$ into spin $15/2$. Their image Gram matrix is $I_{15}/18$, so
\[
 18\Phi(|\psi\rangle\langle\psi|)=I_{16}-|\phi\rangle\langle\phi|.
\]
The complementary vector would be a pure sixth-order AC state on fifteen
qubits. Applying the same fifteen rectangular tensors to that state maps
into the fourteen-dimensional spin-$13/2$ space, contradicting their
orthogonality. Thus $K(17,6)=0$. This proof is independent of the separate
fifth-order exclusion and of the fifteen-qubit fourth-order certificate.

\paragraph{Exact matrix convention.}
In magnetic-number bases, the verified operators are
\[
 \langle j',a|E_{\ell q}|j,m\rangle
 =\sqrt{\frac{2\ell+1}{2j'+1}}
 \langle j,m;\ell,q\mid j',a\rangle,
 \qquad j=17/2,\ j'=15/2.
\]
The verifier checks all Hilbert--Schmidt Gram entries, the raising-operator
covariance equations, absence of every product rank above
$\ell+\ell'$, and the adjoint channel on every diagonal tensor rank from
zero through six. Covariance extends the diagonal checks to every magnetic
component. All these checks use exact SymPy algebraic numbers, with no
floating-point tolerance. The general theorem rests on the proof above,
not extrapolation from a finite computation.
\end{document}
